% Trabalho acadêmico ABNT (NBR 14724) — tese, dissertação, TCC
% Compile: pdflatex main && bibtex main && pdflatex main && pdflatex main

\documentclass[
  12pt,
  oneside,
  a4paper,
  english,
  brazil
]{abntex2}

\usepackage{lmodern}
\usepackage[T1]{fontenc}
\usepackage[utf8]{inputenc}
\usepackage{indentfirst}
\usepackage{graphicx}
\usepackage{microtype}
\usepackage[alf]{abntex2cite}

% --- Metadados (edite aqui) ---
\titulo{Título do trabalho}
\autor{Nome do Autor}
\local{Cidade}
\data{2026}
\orientador{Nome do Orientador}
\instituicao{%
  Nome da Universidade
  \par
  Nome do Programa / Faculdade
}
\tipotrabalho{Dissertação (Mestrado)}
\preambulo{Dissertação apresentada ao Programa X da Universidade Y
como requisito parcial para obtenção do título de Mestre em Z.}

\makeatletter
\hypersetup{
  pdftitle={\@title},
  pdfauthor={\@author},
  pdfsubject={\imprimirpreambulo},
  colorlinks=true,
  linkcolor=black,
  citecolor=black,
  urlcolor=black
}
\makeatother

\setlength{\parindent}{1.3cm}
\setlength{\parskip}{0.2cm}

\begin{document}
\selectlanguage{brazil}
\frenchspacing

% Pré-textuais
\imprimircapa
\imprimirfolhaderosto

\begin{resumo}
Texto do resumo em português.

\vspace{\onelineskip}
\noindent\textbf{Palavras-chave}: palavra1. palavra2. palavra3.
\end{resumo}

\begin{resumo}[Abstract]
\begin{otherlanguage*}{english}
Abstract text in English.

\vspace{\onelineskip}
\noindent\textbf{Keywords}: keyword1. keyword2. keyword3.
\end{otherlanguage*}
\end{resumo}

\tableofcontents*
\textual

\chapter{Introdução}

Texto introdutório. Exemplo de citação: \cite{exemplo2020}.

\section{Objetivos}

Descreva os objetivos do trabalho.

\chapter{Desenvolvimento}

Desenvolva o conteúdo aqui.

\chapter{Conclusão}

Conclusão breve.

\postextual
\bibliography{refs}

\end{document}
